\sectioncentered*{Заключение}
\addcontentsline{toc}{section}{Заключение}

Реализована система автоматического реферирования для варианта~22: русский и немецкий, медицина и критика изобразительного искусства. Классический реферат собирается методом sentence extraction по формулам методички. Рядом строится дерево ключевых слов и семантическая сеть OSTIS: узлы-понятия и дуги-отношения. На выходе есть ссылка на исходный документ, сохранение в файл и печать.

Оценка на четырёх эталонных текстах показывает, что веса отличаются от простого отбора начала: F1 sentence extraction вдвое выше базы. Позиционные функции всё равно опускают поздние тезисы критики. Система применима как учебный стенд: формулы прозрачны, коллекцию можно расширять своими PDF.

Перспектива \mdash учитывать анафору и сжимать выбранные предложения, как допускает методичка для классического реферата.

\begin{thebibliography}{9}
	\bibitem{lab}
	Лабораторная работа <<Автоматическое реферирование документов>>. Кафедра ИИТ БГУИР.
	\bibitem{manning}
	Маннинг К., Рагхаван П., Шютце Х. Введение в информационный поиск. "--- М. : Вильямс, 2011.
	\bibitem{ostis}
	OSTIS. Open Semantic Technology for Intelligent Systems [Электронный ресурс]. "--- Режим доступа: \url{https://ostis-dev.github.io/}.
	\bibitem{fastapi}
	FastAPI documentation [Электронный ресурс]. "--- Режим доступа: \url{https://fastapi.tiangolo.com/}.
	\bibitem{pymorphy}
	Korobov M. Morphological Analyzer and Generator for Russian and Ukrainian Languages // Analysis of Images, Social Networks and Texts. "--- 2015.
\end{thebibliography}
